\section{Moment Estimation}\label{sec:amplitudeEstimation}

\red{Use here amplitude amplification scheme from a different perspective:
    \begin{itemize}
        \item Good and bad states by the restriction to the support of formulas
        \item Good state reflection using the oracle on an disentangled ancilla qubit in the Hadamard basis state $\qstateof{+}$
    \end{itemize}
}

\begin{figure}
    \begin{center}
        \begin{tikzpicture}[scale=0.35, thick] % , baseline = -3.5pt

    \drawsmalltensor{0}{9}{\corelabelsize $\fbasis$}
    \draw (1,9) -- (2,9);
    \drawsmalltensor{3}{9}{\corelabelsize $\paulixsymbol$}
    \draw (4,9) -- (5,9);
    \drawsmalltensor{6}{9}{\corelabelsize $\hgate$}

    %\draw (-1,5) rectangle (1,7);
    \drawsmalltensor{0}{6}{\corelabelsize $\fbasis$}
    \draw (1,6) -- (2,6);
    \drawsmalltensor{0}{2}{\corelabelsize $\fbasis$}
    \draw (1,2) -- (2,2);
    \draw (2,1) rectangle (4,7);
    \drawtextnode{3}{4}{\corelabelsize $\unitarysymbol$}
    \draw (4,6) -- (12,6);
    \draw (4,2) -- (12,2);

    %\draw (2,1) rectangle (4,3);
    %\drawtextnode{3}{2}{\corelabelsize $\hgate$}


    \drawtextnode{-3}{9}{\corelabelsize $\headvariable$}
    \draw[dashed] (-2,7.5) -- (30.5,7.5);
    \drawtextnode{-3}{4}{\corelabelsize $\shortcatvariables$}
    \drawtextnode{0}{4.5}{\corelabelsize $\vdots$}

    \draw[dashed] (7.5,12) -- (7.5,-1) node[below] {\corelabelsize $\qstateofat{0}{\shortcatvariables} \otimes \hplusbasat{\headvariable}$};

    \begin{scope}[shift={(4,0)}]
        \draw (3,9) -- (4,9);
        \draw (4,8) rectangle (7,10);
        \node[anchor=center] (text) at (5.5,9) {\corelabelsize $\qcbencodingof{\exformula}$};
        \draw (5,6) -- (5,8);
        \drawvariabledot{5}{6}
        \draw (6,2) -- (6,8);
        \drawvariabledot{6}{2}
    \end{scope}

    \drawtextnode{9.5}{11}{\corelabelsize $\reflectionof{\badstatesymbol}$}

    \draw[dashed] (11.5,11) -- (11.5,0.5);

    \draw (12,1) rectangle (14,7);
    \drawtextnode{13}{4}{\corelabelsize $\unitaryof{\dagger}$}
    \draw (14,6) -- (15,6);
    \drawsmalltensor{16}{6}{\corelabelsize $\paulixsymbol$}
    \draw (14,2) -- (15,2);
    \drawsmalltensor{16}{2}{\corelabelsize $\paulixsymbol$}
    \draw (17,2) -- (21,2);
    \draw (17,6) -- (18,6);
    \drawsmalltensor{19}{6}{\corelabelsize $\paulizsymbol$}
    \draw (18.5,5) -- (18.5,4);
    \drawvariabledot{18.5}{4}
    \node[rotate=-60] at (18.9,2.75) {$\cdots$};
    \draw (19.5,5) -- (19.5,2);
    \drawvariabledot{19.5}{2}

    \draw (20,6) -- (21,6);
    \drawsmalltensor{22}{6}{\corelabelsize $\paulixsymbol$}
    \draw (23,6) -- (24,6);
    \drawsmalltensor{22}{2}{\corelabelsize $\paulixsymbol$}
    \draw (23,2) -- (24,2);
    \draw (24,1) rectangle (26,7);
    \drawtextnode{25}{4}{\corelabelsize $\unitarysymbol$}
    \draw (26,6) -- (28,6);
    \draw (26,2) -- (28,2);

    \draw (11,9)--(28,9);

    \drawtextnode{19}{11}{\corelabelsize $\reflectionof{0}$}
    \draw[dashed] (26.5,12) -- (26.5,-1) node[below] {\corelabelsize $\qstateofat{1}{\shortcatvariables} \otimes \hplusbasat{\headvariable}$};


\end{tikzpicture}
    \end{center}
    \caption{Amplification scheme of the at the models of a Boolean function $\exformula$ supported coordinates of a quantum state.
    The ancilla variable $\headvariable$ is prepared as an disentagled Hadamard basis vector and stays disentangled during the following reflections.
    The oracle $\qcbencodingof{\formula}$ effectively performs the reflection $\reflectionof{\goodstatesymbol}$ across the good state supported at the models of $\formula$.}
    \label{fig:ancillaGrover}
\end{figure}

Given a distribution $\probwith$ and a propositional formula $\formulaat{\shortcatvariables}$, we now investigate how to use Quantum Phase Estimation to estimate
\begin{align*}
    \meanparam
    = \contraction{\probat{\shortcatvariables},\formulaat{\shortcatvariables}} \, .
\end{align*}

We assume that
\begin{itemize}
    \item a circuit $\unitarysymbol$ preparing a state with measurement distribution $\probwith$, i.e.
    \begin{align*}
        \qstateofat{0}{\ccatvariableof{[\catorder]}{\outsymbol}}
        = \contractionof{\onehotmapofat{\zerotuple{[\catorder]}}{\ccatvariableof{[\catorder]}{\insymbol}},
            \unitaryat{\ccatvariableof{[\catorder]}{\insymbol},\ccatvariableof{[\catorder]}{\outsymbol}}
        }{\ccatvariableof{[\catorder]}{\outsymbol}}
    \end{align*}
    and
    \begin{align*}
        \absof{\qstateofat{0}{\shortcatvariables}}^2
        = \probwith
    \end{align*}
    \item a \computationCircuit{} for $\formulawith$
\end{itemize}


%% FROM AMPLIFICATION
%\subsection{Effective reflection by single \computationCircuits{}}

%\red{Note, that the \computationCircuit{} is not a reflection, but acts as one when preparing the ancilla in the real anti-symmetric state.}
%To start let
%\begin{itemize}
%    \item $\qstateofat{0}{\shortcatvariables}$ has the measurement distribution $\probofat{0}{\shortcatvariables}$
%    \item $\unitarysymbol$ be a unitary preparing the state
%    \item $\qcbencodingof{\exformula}$ be a \computationCircuit{} to a formula $\exformula$
%\end{itemize}

%\subsubsection{Decomposition of the q-sample}
%
%The we split the initial state into vectors supported at the models of $\exformula$ and the complement as
%\begin{align*}
%    \qstateofat{0}{\shortcatvariables}
%    =\qstateofat{\goodstatesymbol}{\shortcatvariables}+\qstateofat{\badstatesymbol}{\shortcatvariables}
%\end{align*}
%where
%\begin{align*}
%    \qstateofat{\goodstatesymbol}{\shortcatvariables}
%    &= \sum_{\shortcatindicesin\wcols\exformulaat{\indexedshortcatvariables}=1} \sqrt{\probat{\indexedshortcatvariables}}
%    \cdot \onehotmapofat{\shortcatindices}{\shortcatvariables} \\
%    \qstateofat{\badstatesymbol}{\shortcatvariables}
%    &= \sum_{\shortcatindicesin\wcols\exformulaat{\indexedshortcatvariables}=0} \sqrt{\probat{\indexedshortcatvariables}}
%    \cdot \onehotmapofat{\shortcatindices}{\shortcatvariables} \, .
%\end{align*}

%We further define for a boolean formula
%\begin{align*}
%    \qstateofat{\probof{0},\exformula}{\shortcatvariables}
%    = \frac{\qstateofat{\goodstatesymbol}{\shortcatvariables}}{\normof{\qstateofat{\goodstatesymbol}{\shortcatvariables}}}
%\end{align*}
%and have consistently for its negation
%\begin{align*}
%    \qstateofat{\probof{0},\lnot\exformula}{\shortcatvariables}
%    = \frac{\qstateofat{\badstatesymbol}{\shortcatvariables}}{\normof{\qstateofat{\badstatesymbol}{\shortcatvariables}}}
%\end{align*}

%We define an angle $\rotanglesymbol$ by
%\begin{align*}
%    \sinof{\frac{\rotanglesymbol}{2}} \coloneqq \normof{\qstateofat{\goodstatesymbol}{\shortcatvariables}}
%\end{align*}
%and have
%\begin{align*}
%    \qstateofat{0}{\shortcatvariables}
%    = \sinof{\frac{\rotanglesymbol}{2}} \qstateofat{\probof{0},\exformula}{\shortcatvariables} +
%    \cosof{\frac{\rotanglesymbol}{2}} \qstateofat{\probof{0},\exformula}{\shortcatvariables} \, .
%\end{align*}
%Further we have that
%\begin{align*}
%    \meanparam^{0}
%    &= \contraction{\probofat{0}{\shortcatvariables},\formulawith} \\
%    &= \sum_{\shortcatindicesin} \absof{\qstateofat{0}{\indexedshortcatvariables}\formulaat{\indexedshortcatvariables}}^2 \\
%    &= \left(\|\qstateofat{\goodstatesymbol}{\shortcatvariables}\|_2\right)^2 \\
%    &= \left(\sinof{\frac{\rotanglesymbol}{2}}\right)^2 \, .
%\end{align*}
%
%We prepare an ancilla qubit in the anti-symmetric state and consider as initial state
%\begin{align*}
%    \qstateofat{0}{\shortcatvariables} \otimes
%    \sqrt{\frac{1}{2}}\cdot\coloredmatrixof{1\\-1}{\avariable} \,.
%\end{align*}
%
%The \computationCircuit{} $\qcbencodingof{\exformula}$ acting on this state prepares a state
%\begin{align*}
%    \left(\qstateofat{\badstatesymbol}{\shortcatvariables}-\qstateofat{\goodstatesymbol}{\shortcatvariables}\right) \otimes
%    \sqrt{\frac{1}{2}}\cdot\coloredmatrixof{1\\-1}{\avariable}  \,.
%\end{align*}
%This is the reflection on the subspace spanned by the one-hot encodings of the models of $\exformula$.



\subsection{Moments }

%We define the moment of the measurement distribution to $\qstateof{0}$ with respect to $\exformula$ as
We have
\begin{align*}
    %\meanparam
    \sum_{\shortcatindices\in[2]^{\catorder} \wcols \formula(\shortcatindices)=1}
    \absof{\qstateofat{0}{\indexedshortcatvariables}}^2
    = \contraction{\probat{\shortcatvariables},\formulaat{\shortcatvariables}}  = \meanparam \, .
\end{align*}
This insight motivates us to define the orthonormal quantum states
\begin{align*}
    \qstateofat{\goodstatesymbol}{\shortcatvariables}
    &= \frac{1}{\sqrt{\meanparam}} \sum_{\shortcatindices\in[2]^{\catorder} \wcols \formula(\shortcatindices)=1}
    \qstateofat{0}{\indexedshortcatvariables} \cdot \onehotmapofat{\shortcatindices}{\shortcatvariables} \\
    \qstateofat{\badstatesymbol}{\shortcatvariables}
    &= \frac{1}{\sqrt{1-\meanparam}} \sum_{\shortcatindices\in[2]^{\catorder} \wcols \formula(\shortcatindices)=0}
    \qstateofat{0}{\indexedshortcatvariables} \cdot \onehotmapofat{\shortcatindices}{\shortcatvariables}
\end{align*}
and split the prepared state into a weighted sum of them as
\begin{align*}
    \qstateofat{0}{\shortcatvariables}
    = \sqrt{\meanparam} \cdot \qstateofat{\goodstatesymbol}{\shortcatvariables}
    + \sqrt{1-\meanparam} \cdot \qstateofat{\badstatesymbol}{\shortcatvariables} \, .
\end{align*}

%% Comparison with quantum rejection sampling
Note that here the mean parameter $\meanparam$ is the analogon to the acceptance probability $\acceptanceprob$ in quantum rejection sampling.
To be more precise, when aiming to sample from the models of $\exformula$ using a rejection sampling scheme starting with the uniform distribution, the acceptance probability would be $\meanparam$.
Instead of amplifying this acceptance probability we here want to measure it to retrieve the mean parameter.


\subsection{Preparation of the effective Grover operator}

The mean parameter can be retrieved from the eigenvalues of an associated (effective) Grover operator as we prepare next.

The Grover operator $\groverofat{\qstateof{0},\formula}{\cheadvariable{\insymbol},\cheadvariable{\outsymbol},\ccatvariableof{[\catorder]}{\insymbol},\ccatvariableof{[\catorder]}{\outsymbol}}$ consists in (see \figref{fig:ancillaGrover})
\begin{itemize}
    \item Effective reflection across $\qstateof{\goodstatesymbol}$ by $\qcbencodingof{\formula}$.
    \item Reflection across $\qstateof{0}$ by $(\unitaryof{\textdagger}\circ\reflectionof{\fbasis}\circ \unitarysymbol)$ where $\reflectionof{\fbasis}$ is the reflection across the ground state $\onehotmapofat{\zerotuple{[\catorder]}}{\shortcatvariables}$ leaving $\headvariable$ invariant.
\end{itemize}


%The Grover operator $\groverofat{\probtensor,\formula}{\shortcatvariables^{\insymbol},\shortcatvariables^{\outsymbol},\avariable^{\insymbol},\avariable^{\outsymbol}}$ amplifying $\formula$
%\begin{align*}
%    (\unitaryof{\textdagger}\circ\reflectionof{0}\circ \unitarysymbol) \otimes \qcbencodingof{\exformula}
%\end{align*}
%acting on the initial state
%\begin{align*}
%    \onehotmapofat{0}{\shortcatvariables} \otimes \sqrt{\frac{1}{2}}\cdot\coloredmatrixof{1\\-1}{\avariable} \, .
%\end{align*}

\subsubsection{Effective reflection by the oracle}

The oracle acts as
\begin{align*}
    \contractionof{
    \qstateofat{0}{\shortcatvariables},\hplusbasat{\cheadvariable{\insymbol}},
    \qcbencodingofat{\exformula}{\cheadvariable{\insymbol},\cheadvariable{\outsymbol}}}{\cheadvariable{\outsymbol},\shortcatvariables}
    =\hplusbasat{\cheadvariable{\outsymbol}}
    \otimes \left(-\sqrt{\meanparam}\qstateof{\goodstatesymbol}+\sqrt{1-\meanparam}\qstateof{\badstatesymbol}\right) \, ,
\end{align*}
and is therefore an effective reflection across $\qstateof{\badstatesymbol}$.
We can exploit the usual sparsity mechanisms to prepare the oracle.



\subsection{Eigenvalues to the formula-amplifying Grover operator}

% Plane projection
We describe the Grover operator by its action on the projection onto the plane
\begin{align*}
    \spanof{\qstateof{\goodstatesymbol},\qstateof{\badstatesymbol}}
\end{align*}
as
\begin{align*}
    \coloredmatrixof{\cosof{\rotanglesymbol} & \sinof{\rotanglesymbol} \\
    -\sinof{\rotanglesymbol} & \cosof{\rotanglesymbol}}{
        \indvariable^{\insymbol},\indvariable^{\outsymbol}
    } \, .
\end{align*}
Here we introduce a boolean variables $\indvariable$ to select the two axis on the plane, and copy it to $\indvariableof{\insymbol},\indvariableof{\outsymbol}$.
The rotation angle is related to $\meanparam$ as
\begin{align*}
    \rotanglesymbol = 2 \cdot \sin^{-1}\left(\sqrt{\meanparam}\right)
\end{align*}

Its eigenvalues are thus
\begin{align*}
    \expof{\pm i\rotanglesymbol}
\end{align*}
to the eigenstates
\begin{align*}
    \frac{1}{2}
    \coloredmatrixof{1\\\pm i}{\indvariable} \, .
\end{align*}
Through measuring the phase $\phasecore\in\{\pm\rotanglesymbol\}$ of one of these eigenvalues we can therefore retrieve the mean by
\begin{align*}
    \meanparam
    = \left(\sinof{\frac{\phasecore}{2}}\right)^2 \, .
\end{align*}

\subsection{Quantum Phase Estimation}

% Generic Quantum Phase Estimation
The phase of the Grover eigenvalues can be estimated using the Quantum Phase Estimation routine \cite{kitaev_quantum_1995, cleve_quantum_1998}.
We prepare auxiliary qubits $\avariables$ and perform for each $\selindexin$ $2^{\selindex}$ repetitions of $\groverof{\probtensor,\exformula}$ controlled on $\avariableof{\selindex}$.
The inverse Fourier transform on the auxiliary qubits is supported on the (2-adic representation of the) phases on the eigenvalues, i.e. $\pm\rotanglesymbol$, and can be detected by a computational basis measure of the ancilla variables.

% For the mean
When measuring the phase register as $\aindexof{[\seldim]}$ we compute a phase
\begin{align*}
    \phasecore
    =2\pi \cdot \sum_{\selindexin} 2^{-(\selindex+1)} \aindexof{\selindex} % 0.\phi_0\phi_1\ldots\phi_{\seldim-1}
\end{align*}
which is either $\alpha$ or $-\alpha$.
We retrieve the mean in both cases by
\begin{align*}
    \meanparam
    = \sin^2\left(\pi \cdot \sum_{\selindexin} 2^{-(\selindex+1)} \aindexof{\selindex}\right) \, .
\end{align*}
